% Source: https://artofproblemsolving.com/community/c3h1869581_summation_notation
% Stable identity: aops-topic-1869581-q01
(a) The terms are $4^0,4^1,4^2,\dots$, so the $k$th term is $4^{k-1}$ and the sum is $\displaystyle\sum_{k=1}^{n}4^{k-1}$. Call it $G$. Then
\[
4G-G=(4+16+\cdots+4^{n})-(1+4+\cdots+4^{n-1})=4^n-1,
\]
so $\displaystyle\sum_{k=1}^{n}4^{k-1}=\boxed{\dfrac{4^n-1}{3}}$. (Check: $n=3$ gives $1+4+16=21=\dfrac{64-1}{3}$.)

(b) The terms increase by $3$, starting at $2$, so the $k$th term is $3k-1$. The last term $302=3k-1$ gives $k=101$. Hence the sum is $\displaystyle\sum_{k=1}^{101}(3k-1)$, and
\[
\sum_{k=1}^{101}(3k-1)=3\cdot\dfrac{101\cdot102}{2}-101=15453-101=\boxed{15352}.
\]
(Equivalently, $101$ terms with average $\dfrac{2+302}{2}=152$ give $101\cdot152=15352$.)
